% ---------------------------------------------------------------------
\chapter{Der Arbeitsablauf}
\label{kap:ablauf}

\begin{lernziele}
\begin{itemize}[nosep]
  \item den Feature-Branch-Workflow kennen, wie ihn die meisten Teams verwenden,
  \item verstehen, \emph{warum} nie direkt auf \cmd{main} gearbeitet wird,
  \item einen vollständigen Arbeitszyklus von Anfang bis Ende durchführen können.
\end{itemize}
\end{lernziele}

\section{Warum nicht einfach auf main arbeiten?}

Technisch kann man auf \cmd{main} committen und direkt pushen. Für ein Wegwerfprojekt geht
das. In der Industrie ist es dagegen fast überall verboten, und zwar aus drei Gründen:

\begin{enumerate}
  \item \textbf{main muss immer funktionieren.} Aus \cmd{main} wird die Software gebaut,
        die Kundinnen und Kunden verwenden. Ein halbfertiger Zwischenstand darf dort nie landen.
  \item \textbf{Vier-Augen-Prinzip.} Jede Änderung wird vor der Übernahme von einer zweiten
        Person geprüft. Das findet Fehler, verbreitet Wissen im Team und verhindert, dass
        jemand allein unbemerkt etwas Falsches einbaut.
  \item \textbf{Automatische Prüfung.} Vor der Übernahme laufen automatisch Tests. Schlägt
        einer fehl, wird die Änderung nicht übernommen.
\end{enumerate}

Diese Prüfungen brauchen einen Ort, an dem die Änderung \emph{wartet}, bevor sie in
\cmd{main} kommt: den eigenen Branch und den Pull Request.

\begin{praxis}
Auf GitHub lässt sich \cmd{main} technisch sperren (\emph{branch protection}). Ein direkter
Push wird dann abgewiesen, und eine Übernahme ist nur über einen Pull Request mit
Freigabe und grünen Tests möglich. In professionellen Teams ist das die Regel. Wer sich
den Ablauf schon bei eigenen Projekten angewöhnt, arbeitet vom ersten Tag an wie im Team.
\end{praxis}

\section{Der Feature-Branch-Workflow}

Der verbreitetste Ablauf heisst \begriff{Feature-Branch-Workflow}: Jede Aufgabe, ob neue
Funktion, Fehlerbehebung oder Übung, erhält einen eigenen kurzlebigen Branch.

\begin{center}
\begin{tikzpicture}[
    schritt/.style={draw=hsluBlau, fill=hsluBlau!6, rounded corners=3pt, thick,
                    minimum width=2.45cm, minimum height=1.05cm, align=center, font=\small},
    weg/.style={-{Stealth[length=2.2mm]}, thick}
  ]
  \node[schritt] (s1) at (0,0)    {aktuellen\\Stand holen};
  \node[schritt] (s2) at (3.0,0)  {Branch\\anlegen};
  \node[schritt] (s3) at (6.0,0)  {arbeiten,\\committen};
  \node[schritt] (s4) at (9.0,0)  {pushen};
  \node[schritt] (s5) at (9.0,-1.8) {Pull Request\\öffnen};
  \node[schritt] (s6) at (6.0,-1.8) {prüfen,\\testen};
  \node[schritt] (s7) at (3.0,-1.8) {in main\\übernehmen};
  \node[schritt] (s8) at (0,-1.8)  {Branch\\löschen};
  \draw[weg] (s1) -- (s2);  \draw[weg] (s2) -- (s3);  \draw[weg] (s3) -- (s4);
  \draw[weg] (s4) -- (s5);  \draw[weg] (s5) -- (s6);  \draw[weg] (s6) -- (s7);
  \draw[weg] (s7) -- (s8);  \draw[weg, dashed, gray] (s8) -- (s1);
\end{tikzpicture}
\end{center}

\section{Ein vollständiger Arbeitszyklus}

Das folgende Beispiel zeigt einen Abend, an dem eine Übung im Modul ADS gelöst wird. Jeder
Befehl wird in Kapitel~\ref{kap:befehle} im Detail erklärt.

\subsection*{1. Aktuellen Stand holen}

\begin{lstlisting}[style=befehle]
git switch main
git pull
\end{lstlisting}

Wer auf einem veralteten Stand beginnt, riskiert später Konflikte mit Änderungen, die in
der Zwischenzeit entstanden sind, zum Beispiel auf einem anderen Gerät.

\subsection*{2. Branch anlegen}

\begin{lstlisting}[style=befehle]
git switch -c ads
\end{lstlisting}

Ab jetzt landen alle Commits auf \cmd{ads}. \cmd{main} bleibt unberührt.

\subsection*{3. Arbeiten und committen}

Dateien bearbeiten, speichern, dann:

\begin{lstlisting}[style=befehle]
git status                             # was hat sich geändert?
git add python                         # Änderungen auswählen
git diff --cached --stat               # Auswahl kontrollieren
git commit -m "ADS U02: Stack mit verketteter Liste implementiert"
\end{lstlisting}

Dieser Schritt darf sich mehrmals wiederholen. Mehrere kleine Commits sind besser als ein
grosser (siehe Kapitel~\ref{kap:commits}).

\subsection*{4. Pushen}

\begin{lstlisting}[style=befehle]
git push -u origin ads
\end{lstlisting}

Jetzt liegt die Arbeit auch auf GitHub und ist gegen den Ausfall des eigenen Computers gesichert.

\subsection*{5. Pull Request öffnen, prüfen, übernehmen}

\begin{lstlisting}[style=befehle]
gh pr create --fill
gh pr merge --merge --delete-branch
\end{lstlisting}

Im Team würde zwischen diesen beiden Befehlen eine Kollegin den Pull Request prüfen und
freigeben. Bei eigenen Projekten übernimmt man diese Rolle selbst und schaut sich die
Änderungen auf GitHub noch einmal in Ruhe an.

\subsection*{6. Abschliessen}

\begin{lstlisting}[style=befehle]
git pull
\end{lstlisting}

\cmd{gh pr merge --delete-branch} hat bereits auf \cmd{main} zurückgewechselt. \cmd{git pull}
holt den neuen Stand von \cmd{main}, der jetzt die eigene Arbeit enthält. Der Kreis ist
geschlossen.

\begin{erfahrung}
\textbf{Merkhilfe für das Ende:} \emph{push} (hochladen) $\rightarrow$ \emph{create}
(Antrag stellen) $\rightarrow$ \emph{merge} (Antrag annehmen) $\rightarrow$ \emph{pull}
(Ergebnis holen). Wer \emph{create} vergisst, erhält beim Mergen die Meldung
„no pull requests found“: Ohne Antrag gibt es nichts anzunehmen.
\end{erfahrung}

\section{Branch-Namen}

In Teams folgen Branch-Namen oft einem Schema, damit man auf einen Blick sieht, worum es
geht, zum Beispiel \cmd{feature/login}, \cmd{fix/absturz-beim-speichern} oder
\cmd{JIRA-1234-export}. Für eigene Projekte genügt ein kurzes Wort ohne Leerzeichen und
Umlaute: \cmd{ads}, \cmd{lial}, \cmd{latex}. Da der Branch nach dem Merge gelöscht wird,
kann derselbe Name immer wieder verwendet werden.

\begin{zusammenfassung}
\begin{itemize}[nosep]
  \item \cmd{main} enthält immer einen funktionierenden Stand und wird nie direkt bearbeitet.
  \item Jede Aufgabe erhält einen eigenen Branch, der per Pull Request übernommen wird.
  \item Zyklus: \cmd{pull} $\rightarrow$ \cmd{switch -c} $\rightarrow$ \cmd{add}/\cmd{commit}
        $\rightarrow$ \cmd{push} $\rightarrow$ \cmd{pr create} $\rightarrow$ \cmd{pr merge}
        $\rightarrow$ \cmd{pull}.
\end{itemize}
\end{zusammenfassung}
